\section{Theoretical Extensions}

\begin{figure*}[t]
\centering
\begin{tikzpicture}[
    phase/.style={rounded rectangle, draw=black!50, minimum width=1.5cm, minimum height=0.55cm, font=\tiny\sffamily, inner sep=2pt},
    emp/.style={phase, fill=cEmpirical!15},
    conj/.style={phase, fill=cConjectured!10, draw=cConjectured!40},
    spec/.style={phase, fill=cSpeculative!10, draw=cSpeculative!40},
    epoch/.style={font=\tiny\sffamily\bfseries, text=black!60},
    arr/.style={-{Stealth[length=3pt]}, thick, black!30},
    every node/.style={align=center}
]

% Empirical phases
\node[emp] (p1) at (0, 0) {Ph.1\\Executor};
\node[emp] (p2) at (2.2, 0) {Ph.2\\+Strategist};
\node[emp] (p3) at (4.4, 0) {Ph.3\\+Cartographer};
\node[emp] (p4) at (6.6, 0) {Ph.4\\+Critic};

\draw[arr] (p1) -- (p2);
\draw[arr] (p2) -- (p3);
\draw[arr] (p3) -- (p4);

% Conjectured phases
\node[conj] (p5) at (8.8, 0) {Ph.5\\Coordinator};
\node[conj] (p6) at (10.6, 0) {Ph.6\\Monitor};
\node[conj] (p7) at (12.4, 0) {Ph.7\\Domains};
\node[conj] (p8) at (14.2, 0) {Ph.8\\Architect};

\draw[arr] (p4) -- (p5);
\draw[arr] (p5) -- (p6);
\draw[arr] (p6) -- (p7);
\draw[arr] (p7) -- (p8);

% Speculative extensions (below, aligned under conjectured)
\node[spec] (pe1) at (8.8, -1.4) {Population\\dynamics};
\node[spec] (pe2) at (10.6, -1.4) {Multi-\\principal};
\node[spec] (pe3) at (12.4, -1.4) {Co-\\evolution};
\node[spec] (pe4) at (14.2, -1.4) {Self-\\reference};

\draw[arr, cConjectured!50] (p8.south) -- (pe4.north);
\draw[arr, cSpeculative!40] (pe1) -- (pe2);
\draw[arr, cSpeculative!40] (pe2) -- (pe3);
\draw[arr, cSpeculative!40] (pe3) -- (pe4);

% Epoch labels
\draw[decorate, decoration={brace, amplitude=4pt}, thick, cEmpirical!50] (-0.8, 0.55) -- (7.4, 0.55);
\node[epoch, text=cEmpirical!70] at (3.3, 1.0) {Empirically observed};

\draw[decorate, decoration={brace, amplitude=4pt}, thick, cConjectured!60] (8.0, 0.55) -- (15.0, 0.55);
\node[epoch, text=cConjectured!80] at (11.5, 1.0) {Theoretically conjectured};

\draw[decorate, decoration={brace, amplitude=4pt, mirror}, thick, cSpeculative!50] (8.0, -1.95) -- (15.0, -1.95);
\node[epoch, text=cSpeculative!70] at (11.5, -2.35) {Speculative extensions};

% Driving force annotation
\node[font=\tiny\sffamily\itshape, text=black!40] at (3.3, -0.7) {Each arrow = bottleneck migration triggers fission};

\end{tikzpicture}
\caption{Roadmap of the framework's progression, with explicit epistemic grading. Blue phases (1 to 4) were directly observed in the Wallfacer experiment. Orange phases (5 to 8) are theoretically conjectured based on the fission principle. Red phases represent speculative extensions to populations, multi-stakeholder settings, and self-referential dynamics. Each transition is driven by bottleneck migration.}
\label{fig:roadmap}
\end{figure*}

The framework developed thus far rests on a single case study (Section 2) and a formal construction (Section 4). In this section, we extend the framework in directions that are increasingly speculative. We present these extensions not as findings but as a theoretical scaffold for future empirical investigation.

\subsection{From Single Systems to Populations}

A natural extension of Phase 8 is to consider multiple system instances running in parallel with heterogeneous configurations. In such a population, a routing mechanism could direct incoming tasks to the best-matched instance, and an Architect could gain organic preference pairs by comparing outcomes across configurations. This parallels the transition from individual learning to population-level selection in evolutionary biology \cite{holland1975}. Boyd and Richerson's \cite{boyd1985} dual-inheritance theory (genetic and cultural) suggests an analogous dual-inheritance model for agentic systems: architectural configuration as the ``genetic'' channel and accumulated knowledge as the ``cultural'' channel.

A shared knowledge layer across instances, which we term a Knowledge Commons, would allow learning to propagate across the population without requiring each instance to rediscover the same patterns. Henrich's \cite{henrich2004} demographic model of cultural evolution predicts that larger populations maintain more complex cultural knowledge; if this analogy holds, agent populations with shared knowledge should exhibit qualitative capability gains relative to isolated instances. However, shared knowledge also introduces the risk of \emph{knowledge pollution}, in which an erroneous belief propagated from one instance infects others, analogous to misinformation cascades in social networks \cite{vosoughi2018}.

\subsection{Multi-Principal Preference Aggregation}

In realistic deployment scenarios, multiple human stakeholders with potentially conflicting preferences interact with the system simultaneously. Arrow's impossibility theorem \cite{arrow1951} establishes that no preference aggregation function simultaneously satisfies a set of desirable axioms (unrestricted domain, Pareto efficiency, independence of irrelevant alternatives, and non-dictatorship). This result suggests that the system should not attempt to compute a single ``optimal'' configuration but should instead maintain and present a Pareto front over the multi-dimensional outcome space \cite{emmerich2006}, allowing human principals to navigate the tradeoff structure themselves.

A further speculative step concerns \emph{emergent initiative}: the possibility that a sufficiently capable system might detect cross-domain patterns invisible to any individual human principal and propose unsolicited changes. This would represent a transition from tool to collaborator, but it immediately raises legitimacy questions about authorization and accountability that connect to ongoing debates about AI alignment \cite{gabriel2020}. We raise this possibility not to endorse it but to identify it as a critical frontier for both technical and normative investigation.

\subsection{Co-evolution and Preference Generation}

If the system's outputs shape its human principals' preferences (by expanding what they consider possible), while the principals' evolving preferences shape the system's behavior, then the two are engaged in co-evolution. This mirrors niche construction in evolutionary biology \cite{odling-smee2003}, in which organisms modify their environment, altering the selection pressures that act on subsequent generations. In such a regime, neither system nor principals converge to a static equilibrium; the dynamic relationship is itself the stable attractor.

At the most speculative extreme, the GP posterior over preference space permits extrapolation into unobserved regions, enabling the system to propose preference hypotheses that its human principals have not yet articulated. Both acceptance and rejection of such proposals constitute informative signal for posterior updates. This idea resonates with Dewey's \cite{dewey1939} pragmatist account of preferences as constructed through inquiry, and with Sen's \cite{sen1977} critique of revealed preference theory. Whether such preference generation is technically feasible, epistemically legitimate, or socially desirable remains entirely open.
